% Generated draft: requires mathematical and source review.
\documentclass[11pt]{article}
\usepackage[margin=25mm]{geometry}
\usepackage{fontspec}
\setmainfont{Latin Modern Roman}
\usepackage{amsmath,amssymb,mathtools,unicode-math}
\setmathfont{Latin Modern Math}
\usepackage{enumitem,needspace,xurl,hyperref}
\setlength{\parindent}{0pt}
\setlength{\parskip}{5pt}
\setlength{\emergencystretch}{3em}
\newcommand{\N}{\mathbb{N}}
\newcommand{\Z}{\mathbb{Z}}
\newcommand{\Q}{\mathbb{Q}}
\newcommand{\R}{\mathbb{R}}
\newcommand{\C}{\mathbb{C}}
\newcommand{\F}{\mathbb{F}}
\newcommand{\A}{\mathbb{A}}
\newcommand{\PP}{\mathbb{P}}
\newcommand{\E}{\mathbb{E}}
\newcommand{\eps}{\varepsilon}
\newcommand{\dd}{\,\mathrm{d}}
\begin{document}
\section*{UnsolvedMath ID: 30004476}
\textbf{Source problem code:} catalogue entry
\subsection{Definitions and mathematical statement}
The supplied record concerns symbols \(p,E,f,g,h,\rho\), a triple-product \(p\)-adic \(L\)-function \(L_p(f,g,h)\), and a twisted \(L\)-function \(s\mapsto L(E,\rho,s)\). It does not specify the admissible data or their hypotheses. In particular, the objects \(f,g,h\), the representation \(\rho\), and the meaning of \(E(\overline{\Q})_\rho\) are not defined in the supplied text. We retain this last notation as the source's designation for the space containing the proposed point; we do not choose an interpretation of its subscript.

A function \(F(s)\) has a simple zero at \(s=1\) if it is holomorphic in a neighborhood of \(1\) and satisfies \(F(1)=0\) and \(F'(1)\ne 0\). Write \(\log_E\) for the formal-group logarithm mentioned in the record. Its normalization, domain, and the choices needed to evaluate it on the proposed point are unspecified.

\textbf{Conjectural target, subject to unresolved source qualifications.} For the data and hypotheses intended by the source, there exists a canonically specified multiple \(L_p^{\mathrm{can}}(f,g,h)\) of \(L_p(f,g,h)\), together with a point described by the source as a rational point
\[
P_\rho\in E(\overline{\Q})_\rho,
\]
such that
\[
L_p^{\mathrm{can}}(f,g,h)(2,1,1)=\log_E(P_\rho)
\]
and
\[
P_\rho\text{ is nontrivial}
\quad\Longleftrightarrow\quad
\bigl(s\mapsto L(E,\rho,s)\bigr)\text{ has a simple zero at }s=1.
\]
Here \((2,1,1)\) is the evaluation point appearing in the source; its parameter space is not supplied. The permitted multiplier and its canonical characterization are also unspecified. The terms ``rational point'' and ``nontrivial'' are retained without identifying a field of rationality or choosing between nonzero and non-torsion interpretations. Thus the displayed target records the supplied conjecture, but the record does not support a fully quantified formulation over a defined class of inputs.
\subsection{Short English statement}
The conjecture asks for a canonical normalization of a triple-product \(p\)-adic \(L\)-function whose value at \((2,1,1)\) is the formal-group logarithm of a point \(P_\rho\). That point should be nontrivial exactly when \(L(E,\rho,s)\) has a simple zero at \(s=1\). The admissible data, normalization, point space, and meaning of nontriviality require clarification from the source.
\subsection{Sources}
\begin{itemize}
\item Supplied local UnsolvedMath registry record, permanent ID \(30004476\), exact source problem code \texttt{catalogue entry}, entitled ``Canonical Triple-Product \(p\)-Adic \(L\)-Values and Rational Points,'' in the Oberwolfach Reports Open Problems collection. The record gives proposed year 2020 and names no individual proposer.
\item Supplied source citation: \emph{Arithmetic Geometry (2021)}, Oberwolfach Reports, DOI 10.4171/owr/2020/20, pp. 1023--1082. Original OWR report: \url{https://doi.org/10.4171/owr/2020/20}. The supplied assessment identifies the target as Conjecture 5; no more precise location is supplied.
\item Supplied literature source: ``Darmon--Rotger diagonal cycles II,'' \url{https://web.mat.upc.edu/victor.rotger/docs/DR2.pdf}.
\item Supplied literature source: ``p-adic families of diagonal cycles,'' \url{https://arxiv.org/abs/2207.01314}.
\item Literature checked at: \texttt{2026-08-22}. This is supplied dated curation, not an independently verified assessment. The record labels the problem partially solved and states that the Conjecture 5 value and nontriviality criterion remain unproved in general, while Darmon--Rotger establish new analytic-rank-zero cases and conditional higher-rank Selmer consequences. Its background characterizes this as literature triage derived from the supplied OWR curation and requests independent verification of time-sensitive status and cited claims.
\item The registry assigns the statement status \texttt{corrected\_verified} and reports a check against the official Oberwolfach report, asserting preservation of its hypotheses, notation, and target. This is a supplied verification claim; no independent source consultation or mathematical verification is asserted here.
\end{itemize}
\end{document}
